\section{Exact feasibility and the resistive model}
\label{sec:foundations}

In a resistive network, current depends linearly on bus voltages, but power is
the product of voltage and current. This distinction makes exact feasibility a
quadratic problem even when the network has no discrete controls. We study its
complexity with interval restrictions on both voltage and net power injection.
The model concerns direct-current electrical networks; it is distinct from
the linear ``DC approximation'' of an alternating-current network.

\subsection{Input model and complexity convention}

Let $G=(N,E)$ be a finite simple undirected graph. Each line $\{i,j\}\in E$
has conductance $g_{ij}=g_{ji}\in\Q_{>0}$. Bus $i$ has voltage bounds
$0<\ell_i\le u_i$ and injection bounds $p_i^-\le p_i^+$, all rational and
finite. Its net injection at voltage vector $v\in\R^N$ is
\begin{equation}
 \power_i(v)=v_i\sum_{j:\{i,j\}\in E}g_{ij}(v_i-v_j).
 \label{eq:rpf}
\end{equation}
An injection is positive for net supply to the network and negative for net
withdrawal. An isolated bus has injection zero. The decision problem $\RPF$
asks whether
\begin{equation}
 \ell_i\le v_i\le u_i,\qquad
 p_i^-\le\power_i(v)\le p_i^+\qquad(i\in N)
 \label{eq:rpf-bounds}
\end{equation}
has a solution. Singleton intervals are allowed. We impose neither additional
line-flow limits nor a requirement that an injection be fixed whenever a
voltage is variable. Thus the input permits voltage-pinned buses with fixed
injections, variable-voltage buses with fixed injections, and buses with both
quantities in intervals. These choices are material to the reduction.

The graph and its rational data are encoded in binary. The existential
theory of the reals, denoted ETR, asks whether an existentially quantified
Boolean combination of polynomial equalities and inequalities with integer
coefficients is true. The class $\ER$ consists of decision problems reducible
to ETR by polynomial-time many-one reductions. Rational coefficients are
equivalent: each atomic constraint can have its denominators cleared using
a positive integer of polynomial bit length. We use the Turing model of
computation and the standard inclusions
\begin{equation}
 \NP\subseteq\ER\subseteq\PSPACE.
 \label{eq:complexity-inclusions}
\end{equation}
See \citet[Section 1.1]{AbrahamsenMiltzow2019} for these conventions and their
background. Exact feasibility quantifies over real solutions, not over decimal
approximations. An irrational solution alone does not rule out membership in
$\NP$: a certificate need not list rational bus voltages.

\begin{proposition}
\label{prop:rpf-membership}
The problem $\RPF$ belongs to $\ER$.
\end{proposition}
\begin{proof}
Equations~\eqref{eq:rpf}--\eqref{eq:rpf-bounds} form a conjunction of rational
polynomial inequalities of degree at most two in $|N|$ real variables.
Writing each incident-line term explicitly uses $O(|N|+|E|)$ terms and
polynomially many bits. Clearing denominators gives an ETR instance.
\end{proof}

The model respects the usual dissipation identity:
\begin{equation}
 \sum_{i\in N}\power_i(v)
 =\sum_{\{i,j\}\in E}g_{ij}(v_i-v_j)^2\ge0.
 \label{eq:dissipation}
\end{equation}
Indeed, the two terms of a line sum to $g_{ij}(v_i-v_j)^2$.
The injection intervals, rather than a separate balance equation with zero
total injection, account for these losses. In particular, permitting signed
injections does not remove physical power balance.

\subsection{The bounded arithmetic source problem}

An instance of $\ETRINV$ consists of named variables $x_1,\ldots,x_n$,
each restricted to $[1/2,2]$, and a conjunction of equations of the forms
\begin{equation}
 x=1,\qquad x+y=z,\qquad xy=1.
 \label{eq:etr-inv}
\end{equation}
Variable names in an equation may coincide. The interval restrictions are
part of feasibility, not a promise that all real solutions already lie there.
Abrahamsen, Adamaszek, and Miltzow prove this problem $\ER$-complete
\citep[Definition 5 and Theorem 7 of the STOC/arXiv version]{AbrahamsenEtAl2018}.
The journal version is \citet{AbrahamsenEtAl2022}.
Replacing each equation $x=1$ by $x x=1$ preserves the solution set, since
$x>0$. Hence we use only addition and inversion equations below.
Let $S_\Phi\subseteq[1/2,2]^n$ denote the solution set of the normalized
instance $\Phi$, and let $m$ be its number of equations.
We retain all named variables, including those occurring in no equation.
